% Working notes -- vertex cover on nearly complete graphs, as a chygraph
% problem.  Temporary: chygraph/vc/.  Numbers come from the scripts in this
% folder (ensemble.py, sbm_sizes.py, chybp.py, compare.py; results/).
\documentclass[11pt]{article}
\usepackage{amsmath,amssymb,booktabs,graphicx,hyperref}
\usepackage[margin=1.1in]{geometry}
\newcommand{\ave}[1]{\langle #1\rangle}
\newcommand{\VC}{\mathrm{VC}}
\title{Vertex cover on nearly complete graphs:\\ dense complexes, tree-like propagation}
\author{working notes}
\date{\today}

\begin{document}
\maketitle

\section{The question}

Figure 11.5 of the book runs vertex cover along the assortativity axis
$e_{dd'}=q_d[r\delta_{dd'}+(1-r)q_{d'}]$ at fixed $p_d\sim d^{-\tau}$. Left of
$r_c$ the graph is tree-like and replica symmetry is exact; right of $r_c$ the
fixed point of the hard-field map is not an answer. The standard move is
replica-symmetry breaking. The question here is whether the $r\to1$ end has a
better description than that, since at $r=1$ the ensemble is a union of
degree-class blocks, and the starved high-degree classes are nearly complete
graphs. What follows is not a replica ansatz. It is a decomposition.

\section{The complete graph and its neighbourhood}

\paragraph{$K_m$.} An uncovered set is independent and the largest independent
set of $K_m$ is a vertex, so the minimum cover is $m-1$, attained by the $m$
covers that leave one vertex out. Any two differ in two vertices, all pairwise
overlaps equal $1-2/m$: the ground states form a simplex, which is the
replica-symmetric overlap structure. The partition function at chemical
potential $\mu$ per uncovered vertex is $Z=1+m e^{\mu}$ (times the cover
prefactor), entropy $\ln m$, non-extensive. Nothing here calls for RSB.

What fails on $K_m$ is the Bethe approximation, not the ansatz. The hard-field
cavity map of the book, Eq.~(11.11) with one class of excess degree $m-1$, is
$\pi=(1-\pi)^{m-1}$, giving $\pi\simeq\ln m/m$ and an uncovered set of
$\simeq(\ln m)^2/2$ vertices instead of one (nine at $m=100$, nineteen at
$m=1000$), with a Jacobian $|J|=(m-1)(1-\pi)^{m-2}\sim\ln m$ that grows without
bound. The cavity treats the $m-1$ neighbours as conditionally independent; on
a clique they are maximally dependent.

\paragraph{$K_m-H$.} Complementation turns the problem inside out: an
independent set of $K_m-H$ is a clique of $H$,
\begin{equation}
  Z_{K_m-H}(\lambda)=\sum_{C\ \text{clique of }H}\lambda^{|C|}
  =1+m\lambda+|E(H)|\lambda^2+T(H)\lambda^3+\dots
  \label{eq:cliquepoly}
\end{equation}
truncating at $\omega(H)$. For a sparse random $H$, $\omega\le3$ w.h.p. and
the coefficients are the vertex, edge and triangle counts of $H$. The dense
side of vertex cover therefore has its own exact expansion, in the missing
links rather than the present ones, and it is a finite one: the cover is
$m-\omega(H)$, exposed set a maximum clique of $H$, degeneracy $N_\omega(H)$.
Bethe expands around the empty graph in present links; this expands around the
complete graph in missing links. The reason the dense side is trivial is that
the maximum clique of a sparse graph is bounded.

\section{The bracket}

Let $\{b\}$ be vertex-disjoint dense blocks of $G$ (cliques, or $K_{m_b}-H_b$
with $H_b$ sparse) and $R$ the rest of the graph, the blocks deleted. Vertex
cover is superadditive over vertex-disjoint induced subgraphs, and covering
every block entirely is a valid cover, so
\begin{equation}
  \VC(R)+\sum_b\bigl(m_b-\omega_b\bigr)\ \le\ \VC(G)\ \le\ \VC(R)+\sum_b m_b .
  \label{eq:bracket}
\end{equation}
The width is $\sum_b\omega_b$, one vertex per clique, and both ends are
computations on $R$, which is sparse. Two regimes:
\begin{enumerate}
\item Few or small blocks: the width is $O(\#\text{blocks})$ and the
  ambiguity is resolved exactly --- each block exposes one clique of $H_b$,
  chosen among $N_\omega$ candidates, the exposed vertices force their
  external neighbours into the cover, and this is a hard-core problem on the
  block graph with a finite alphabet per block. It is the chygraph cavity with
  the block as a complex: Eq.~(11.3) of the book with the $\max$ taken over
  cliques of $H_b$ rather than over members.
\item Large blocks: the width is $o(\sum_b m_b)$, and the lower end is
  attained as soon as a block has a candidate exposure whose external
  neighbours are already covered, expected number $N_\omega\,x^{\omega
  d_{\rm ext}}\gg1$. The block then costs exactly $m_b-\omega_b$ and imposes
  nothing on the outside: delete it, cover all but $\omega_b$, solve $R$.
\end{enumerate}
The tree-like approach to the $r\to1$ limit is therefore: remove the dense
blocks, run the ordinary cavity on what is left, add
$\sum_b(m_b-\omega_b)$. The dense part contributes a number, not a field.
A minimum cover contains at least $m-1$ vertices of every $m$-clique, so the
dense part is forced and only its complement is a statistical-mechanics
problem.

\paragraph{Where one vertex per block is not negligible.} For the cover size it
is. For the core it is not: leaf removal never enters a clique of size $\ge3$,
so every block sits in the core whole, however large. The lower panel of
Fig.~11.5 is an $O(1)$-per-block-sensitive diagnostic and $x_c$ is not; near
$r=1$ the two can disagree without either being wrong.

\paragraph{Where the instability of Eq.~(11.11) actually lives.} The Jacobian
along the $r$ axis is diagonal plus rank one, and the unstable eigenvalue
solves $1=(1-r)\sum_d q_d a_d/(\Lambda-r a_d)$ with
$a_d=d(1-\pi_d)^{d-1}$. Weighting each class by its term in that sum (cutoff
$10^3$; reproduces $r_c=0.70$, $0.78$):
\begin{center}
\begin{tabular}{lc}
\toprule
 & classes carrying the instability (quartiles, 90\%, 99\%)\\
\midrule
$\tau=2.5$, $r=r_c=0.70$ & $k=3,\ 5,\ 6,\ 9,\ 18$\\
$\tau=3.0$, $r=r_c=0.78$ & $k=3,\ 4,\ 6,\ 7,\ 13$\\
$\tau=2.5$, $r=0.99$ & $k=172,\ 174,\ 175,\ 178,\ 205$\\
\bottomrule
\end{tabular}
\end{center}
At $r_c$ the instability is in $k=3$--$9$: at $n=2\times10^5$ these are
random regular graphs of $10^3$--$10^4$ vertices, genuinely sparse, and the
transition is theirs. As $r\to1$ the instability migrates to high $k$; by
$r=0.99$ it sits at $k\approx170$--$200$, where the class saturates
($np_k\lesssim k$, i.e. $k\gtrsim(n/\zeta(\tau))^{1/(\tau+1)}\approx30$ at
$\tau=2.5$, $\approx20$ at $\tau=3$) or is empty. The right-hand end of the
dashed curve is the map describing classes that in the graphs are cliques or
absent. That end is not ``unstable, go to RSB''; it is outside the map's
domain, and Eq.~\eqref{eq:bracket} is its answer.

\section{The chygraph version: communities as complexes}

The proposal to test: take a graph, find its communities with the nested
degree-corrected SBM, make each community a complex, solve exactly inside the
complex and propagate tree-like between complexes. Two questions.

\subsection{Are the communities small and $n$-independent here?}

In the growth models and the real networks of the community-persistence
programme the bottom-level SBM communities have of order $20$--$100$ vertices
and that number does not grow with $n$. On the degree-correlated ensemble it
does not hold. Nested DC-SBM (\texttt{sbm\_sizes.py}, $\tau=2.5$, one seed):
\begin{center}
\small
\begin{tabular}{rrrrrrrl}
\toprule
$r$ & $n$ & blocks & node-median size & max & nodes in blocks $\le50$ & internal edges & what the blocks are\\
\midrule
0.0 & $2\times10^3$--$5\times10^4$ & 1 & $n$ & $n$ & 0 & 1.00 & nothing to find\\
0.5 & $2\times10^3$--$2\times10^4$ & 1--2 & $n$ & $n$ & 0 & 0.93--1.00 & at most one hub block\\
0.5 & $5\times10^4$ & 4 & $n$ & 48559 & 0 & 0.86 & two hub blocks ($\ave k\approx7,10$, sparse)\\
0.8 & $2\times10^3$ & 5 & 1387 & 1387 & 0.01 & 0.95 & degree classes, loosely\\
0.8 & $5\times10^4$ & 34 & 67 & 11586 & 0.00 & 0.61 & degree classes, loosely\\
1.0 & $2\times10^3$ & 8 & 1522 & 1522 & 0.04 & 1.00 & degree classes exactly (sd $k\approx0$)\\
1.0 & $2\times10^4$ & 23 & 37 & 5290 & 0.01 & 0.65 & degree classes exactly\\
1.0 & $5\times10^4$ & 30 & 282 & 30961 & 0.01 & 0.99 & degree classes exactly\\
\bottomrule
\end{tabular}
\end{center}
Three readings. At $r\le0.5$ the degree-corrected model finds one block:
the configuration model \emph{is} its null, and there are no communities to
find. At $r=1$ the blocks are the degree classes to the digit
($\ave{k}=1.00,2.00,2.98,3.96,\dots$ with standard deviation $\lesssim0.3$),
so their sizes are $np_k$ and grow linearly with $n$; the largest block is
the $k=1$ class, a perfect matching on half the graph. Between, at $r=0.8$,
the blocks are degree classes blurred by the $(1-r)$ random links. The only
small, $n$-stable blocks are the saturated classes, of size $\approx k+1$ for
$k\gtrsim(n/\zeta)^{1/(\tau+1)}$, which hold a few per cent of the vertices
at most. So the community scale of this ensemble is not $50$: it is the class
size $np_k$, and the statement from the growth models does not transfer.

\subsection{Exact interior, tree-like exterior}

\texttt{chybp.py} implements the scheme regardless: regions are the SBM
blocks of at most $s_{\max}$ vertices (singletons otherwise, so $s_{\max}=0$
is the node-level cavity), the interior is a weighted vertex cover solved by
branch and bound, the exterior is zero-temperature min-sum with the exact
region-to-vertex message, and the min-sum Bethe energy
$E=\sum_R\varepsilon_R+\sum_e\varepsilon_e-\sum_i(d_i-1)\varepsilon_i$ is
exact when the region graph is a tree (checked against brute force on 200
random tree-glued instances). A valid cover is read off by decimation and
compared with the exact optimum from an integer program.

% RESULTS TABLE: filled from results/compare.json
\input{results/compare_table.tex}

\section{Line graphs: dense-block chygraphs that are RS-exact at every density}

If every vertex is in at most two cliques (chy-degree $\le2$), $G$ is a line
graph (Krausz) and an independent set is a set of vertices no two sharing a
clique, i.e.\ a matching in the clique graph $D$ (nodes $=$ complexes, edges
$=$ vertices, node degree $=$ cardinality):
\begin{equation}
  \VC(G)=|V|-\nu(D).
\end{equation}
Maximum matching is polynomial (Edmonds), has no phase transition on any
graph at any fugacity (Heilmann--Lieb), is RS-exact under the cavity method
(Zdeborov\'a--M\'ezard 2006), and its ground state is Karp--Sipser (Poisson) or
Bohman--Frieze 2011 (arbitrary degree sequence, here the cardinality
distribution). Cardinality is irrelevant to hardness: it is only the dual
degree. Chy-degree is the dual hyperedge size, and at $\ge3$ the dual is
hypergraph matching, where NP-hardness and presumably RSB live. This is the
dual of the book's Sec.~11.6: vertex-level leaf removal never certifies at
$c\ge3$, but for chy-degree $\le2$ a certificate exists, Karp--Sipser on $D$.

\section{Literature}

Replica theory near the complete graph is the large-connectivity expansion
around SK (Kanter--Sompolinsky 1987; de Dominicis--Goldschmidt 1989;
M\'ezard--Parisi 2001, $c\to\infty$ of the Bethe spin glass), and there the
correct ansatz becomes Parisi full RSB, not something simpler. For vertex
cover the dense limit is degenerate: on $G(n,p)$ at fixed $p$,
$\alpha=O(\log n)$, and on $K_n-H$, $\VC=n-\omega(H)$. The nontrivial dense
results are the large-$c$ asymptotics (Frieze 1990, matched at leading order
by the Weigt--Hartmann RS formula), the planted-clique analysis
(Deshpande--Montanari 2015, AMP as the cavity on $K_n$) and the algorithmic
gap (Karp's $\log_2n$ vs $2\log_2n$; Gamarnik--Sudan's overlap gap property).
On the object that dominates $x_c(r=1)$, the random regular graph, the
answer is 1RSB (Zhou 2003; Weigt--Zhou 2006; Barbier et al.\ 2013) and proven
so at large degree (Ding--Sly--Sun 2016).

\end{document}
